\section{Historia de la evolución de las redes sociales: de las plataformas de contenido a la socialización AI Native}

La trayectoria de Tristan va de una plataforma de contenido de audio, la socialización entre desconocidos, los juegos de citas y la compañía con IA, hasta Elys. Este recorrido puede verse como una línea de evolución de los productos sociales: primero, la distribución de contenido; después, el emparejamiento entre personas; luego, personajes virtuales que aportan valor emocional, y por último, dobles digitales de IA que ayudan de forma proactiva a conectar a personas reales.

\subsection{Trayectoria de evolución}

\noindent\begin{longtable}{p{0.22\textwidth}p{0.34\textwidth}p{0.35\textwidth}}
\toprule
Etapa & Mecanismo central & Limitaciones \\
\midrule
Plataforma de contenido & Distribución personalizada de audio/contenido & Conecta a las personas con el contenido, no a las personas entre sí. \\
Socialización entre desconocidos & Etiquetas de baja dimensión y emparejamiento inmediato & Romper el hielo es costoso; la calidad de las relaciones es inestable. \\
Juegos de citas/compañía con IA & El contenido o la IA aporta valor emocional & El interlocutor puede no ser una persona real. \\
Red social de IA & El doble digital, con su context, conecta de forma proactiva a personas reales & Privacidad, confianza, proporción de personas reales y dificultad del arranque en frío. \\
\bottomrule
\end{longtable}

\begin{knowledgebox}{Por qué Elys no surgió de la nada}
Elys hereda la experiencia de la recomendación de contenido, la socialización entre desconocidos, los juegos de citas y la compañía con IA, pero intenta llevar el objetivo de «consumir contenido/compañía» a «conectar a personas reales».
\end{knowledgebox}

\subsection{Resumen de la sección}

La socialización AI Native no aparece de la nada: es la confluencia de tres líneas, las plataformas de contenido, el emparejamiento social y la compañía con IA. La novedad de Elys consiste en que el doble digital convierte de forma proactiva el context en conexiones entre personas reales.
